\documentclass[11pt]{article}

\usepackage[margin=1in]{geometry}
\usepackage{amsmath,amssymb,amsthm}
\usepackage{mathtools}
\usepackage{booktabs}
\usepackage{array}
\usepackage{hyperref}
\usepackage{url}
\usepackage{graphicx}
\usepackage{xcolor}
\usepackage{listings}
\usepackage{enumitem}

\lstset{basicstyle=\ttfamily\small, columns=fullflexible, breaklines=true,
        frame=single, xleftmargin=\parindent}
\hypersetup{colorlinks=true, linkcolor=blue!50!black, urlcolor=blue!50!black,
            citecolor=blue!50!black}

\title{Silicon-Aware Local RTRL on a Hardware-Primitive Grid Lattice}

\author{Thomas Pluck}
\date{\today}

\begin{document}

\maketitle

\begin{abstract}
A recurrent architecture built from one analog-crossbar primitive
(Gm--C integrator $+$ $\tanh$ $+$ learnable bias) replicated over a
rectangular $N{\times}M$ grid. Each cell has two weighted inputs
(from-left and from-above) and one output routed to both its right
and below neighbours: the left shoreline takes external input, the
bottom shoreline produces block output, and the top/right shorelines
are nulled. Four learnable scalars per cell. The encoder (left
shoreline) and decoder (bottom shoreline) carry no learned parameters.
We train the grid on row-wise sequential MNIST with a local SnAP-$1$
rule (two descendant cross-traces per cell); per-parameter grid
gradients match BPTT at float-$64$ $\epsilon$ on a single forward step,
confirming numerical correctness. Training numbers are in
Section~\ref{sec:mnist}.
\end{abstract}

\section{Introduction and contribution}

\paragraph{Setting.}
Analog neuromorphic crossbars can implement a Gm--C low-pass integrator
with a $\tanh$-shaped output stage cheaply and uniformly. The design
question is not whether such a substrate can be made to learn at all,
but what \emph{architecture} and \emph{rule} fit inside its constraints:
no backpropagation-through-time (no global backward pass), no
per-parameter adaptive optimiser state, bounded per-cell analog state,
finite arithmetic precision, fixed input wiring, and spatial locality of
the learning signal. We take those constraints as the scope of the
problem rather than axes to apologise along.

\paragraph{Contributions.}
\begin{enumerate}[leftmargin=*]
\item \textbf{A grid primitive with two-in/two-out connectivity.} An
$N{\times}M$ rectangular grid of cells where each cell receives inputs
from its left and above neighbours and drives output to its right and
below neighbours. Left shoreline is external input, bottom shoreline
is block output, remaining shorelines are nulled. Four learnable
scalars per cell $(\log\tau, w_\text{L}, w_\text{T}, b)$.
\item \textbf{A local rule with per-descendant cross-traces.} SnAP-$1$
on this topology tracks one self-trace per learnable parameter plus
one cross-trace per descendant (2 per cell). Single-step gradients
match BPTT at float-64 $\epsilon$. Full-sequence gradients diverge
with $T$ because each cell has two descendant paths that SnAP-$1$ only
partially captures---a harder truncation regime than the tree variant,
which had one descendant per cell.
\item \textbf{A silicon-aware characterisation.} Within the scope the
hardware imposes, we map six design-relevant axes: weight/state/trace
precision, SnAP-$k$ accuracy vs.\ state budget, timescale range,
topology $(T, K, B, s)$ at matched cell count, capacity scaling, and
transfer to an RL task (MinAtar) and a neuromorphic benchmark (SHD).
The output is a design-space picture, not a leaderboard.
\end{enumerate}

\paragraph{Scope.}
The left and bottom shorelines of the grid are used as input and
output respectively, carrying no learned parameters; when input
dimension matches grid height the encoder is pure identity wiring.
The only trainable parameters are the four per-cell scalars.

\paragraph{Relation to prior work.}
The RTRL taxonomy has been mapped by Marschall, Cho and Savin
\cite{marschall2020}; sparse and stochastic approximations include
UORO~\cite{tallec2018}, KF-RTRL~\cite{mujika2018}, and
SnAP-$k$~\cite{menick2021}. RFLO~\cite{murray2019} and e-prop
\cite{bellec2020} are the closest local-rule neighbours. OSTL
\cite{bohnstingl2023} and Zucchet \emph{et al.}~\cite{zucchet2023} both
exploit structural sparsity for exact online learning in specific
architecture classes; our tree variant is a generalisation. On the
\emph{dendritic} side, models combining multi-compartment neurons with local
learning form a growing literature. The nearest work is Yin \emph{et al.}
\cite{yin2024}: they apply e-prop (equivalent to SnAP-$1$) to a
two-compartment spiking dendritic neuron, on SHD. Differences from the
present work: (i) we generalise from SnAP-$1$ to arbitrary $k$ on an
ancestor chain and make the structural-sparsity argument explicit; (ii)
we work with arbitrary tree depth and branching $(T, K, B, s)$ rather than
a fixed two-compartment topology; (iii) the Gm--C$+\tanh$ primitive is
fully differentiable, whereas LIF-based work requires surrogate
gradients---a genuine simplification for an exact-RTRL argument and for
analog mapping. Zhang \emph{et al.}'s TC-LIF \cite{zhang2024tclif},
Huang \emph{et al.}'s DendSN \cite{huang2024dendsn}, and Feng \emph{et
al.}'s TS-LIF \cite{feng2025tslif} are parallel dual-compartment
proposals. Biologically motivated multi-compartment learning
rules include Sacramento \emph{et al.}~\cite{sacramento2018},
Guerguiev \emph{et al.}~\cite{guerguiev2017}, and
Payeur \emph{et al.}~\cite{payeur2021}. On the ML-side dendrite literature:
Beniaguev \emph{et al.}~\cite{beniaguev2021} show single cortical neurons
match small deep ANNs; Iyer \emph{et al.}~\cite{iyer2022} use active
dendrites for multi-task generalisation; Poirazi and
Papoutsi~\cite{poirazi2020} review the field. On the analog hardware side,
Cartiglia \emph{et al.}~\cite{cartiglia2022} and related work from the
Indiveri group including Maryada \emph{et al.}~\cite{maryada2025} prototype
dendritic learning in silicon; the review of
Acharya \emph{et al.}~\cite{acharya2021} covers the broader landscape.

What is novel here is the combination of: (i) a rectangular-grid
primitive with 2-in/2-out connectivity and native shoreline I/O,
giving the same filter+$\tanh$+bias cell a very different sensitivity
structure from the tree variant \cite{tree-neurons}; (ii) a SnAP-$1$
specialisation that tracks one cross-trace per immediate descendant
(2 per cell) rather than a single ancestor chain; (iii) empirical
quantification of the multi-hop-descendant truncation cost that
SnAP-$1$ pays on this topology, which is significantly higher than on
the tree. What is \emph{not} novel is the SnAP-$k$ framework, the
per-group-LR optimiser story, or the fixed-shoreline-encoder approach.

\paragraph{Code.} \url{https://github.com/ThomasPluck/grid-neurons}

\section{Architecture}
\label{sec:arch}

\subsection{The canonical cell}

Every grid cell is the same first-order integrator followed by $\tanh$.
Cell $(i, j)$ at position row $i$, column $j$ receives two inputs:
one from its left neighbour $y_{i, j-1}$ and one from its top neighbour
$y_{i-1, j}$, each weighted independently:
\begin{align}
\mathrm{drive}_{ij}(t) &= w_{ij}^{\text{L}}\, y_{i,j-1}(t) + w_{ij}^{\text{T}}\, y_{i-1,j}(t), \\
\tau_{ij} \frac{\mathrm{d}s_{ij}}{\mathrm{d}t} &= -s_{ij} + \mathrm{drive}_{ij}(t), \\
y_{ij}(t) &= \tanh\!\big(s_{ij}(t) + b_{ij}\big).
\end{align}
Zero-order-hold discretisation:
\begin{equation}
a_{ij} = \mathrm{e}^{-\Delta t/\tau_{ij}}, \quad c_{ij} = 1 - a_{ij},\qquad
s_{ij}(t{+}1) = a_{ij}\, s_{ij}(t) + c_{ij} \cdot \mathrm{drive}_{ij}(t{+}1).
\label{eq:discrete}
\end{equation}
Each cell carries four learnable scalars,
$\theta_{ij} = \{\log \tau_{ij},\, w_{ij}^{\text{L}},\, w_{ij}^{\text{T}},\, b_{ij}\}$.
Output $y_{ij} \in (-1, 1)$ is bounded, so no normalisation is required.

\subsection{Topology: rectangular grid with shoreline I/O}

A grid block is a rectangle of cells indexed $(i, j)$ with
$i \in \{0, \ldots, N{-}1\}$, $j \in \{0, \ldots, M{-}1\}$. Connectivity
is fixed and feed-forward within a timestep:
\begin{itemize}[leftmargin=*]
\item Each cell has \textbf{two inputs} from its left and top neighbours,
and \textbf{two output edges} routing $y_{ij}$ to its right and below
neighbours.
\item \textbf{Left shoreline} $(j{=}0)$: external input. Each cell in
column $0$ receives one external channel in place of a left neighbour.
\item \textbf{Bottom shoreline} $(i{=}N{-}1)$: block output. Each cell in
row $N{-}1$ is one of the $M$ output values.
\item \textbf{Top shoreline} $(i{=}0)$: the missing top neighbour is
treated as zero. \textbf{Right shoreline} $(j{=}M{-}1)$: the outgoing
right-edge is dropped.
\end{itemize}
Forward traversal is raster order (left-to-right within a row, then
top-to-bottom): each cell's inputs come from cells already updated this
step. Backward traversal is the reverse. The grid has $NM$ cells and
$\sim 2NM$ internal edges, so the forward pass is $O(NM)$ per timestep.

\subsection{Encoder and decoder: native I/O}
\label{sec:encdec}

A block is one grid; it holds every trainable parameter. Encoder and
decoder are defined by the grid's shoreline geometry and carry no
learned parameters.

\paragraph{Encoder.} The left shoreline directly accepts external input:
cell $(i, 0)$'s left-input at each timestep is external channel $x_i(t)$.
For tasks where $N_{\text{ext}} = N$ (input dim $=$ grid height) this is
literally an identity wiring; when $N_{\text{ext}} \ne N$ an optional
fixed sparse-random projection $W_{\text{in}} \in \mathbb{R}^{N \times
N_{\text{ext}}}$ maps external channels to left-shoreline drives. The
projection is drawn once at initialisation and not trained.

\paragraph{Decoder.} The bottom shoreline is the $M$-dimensional block
output. Cell $(N{-}1, j)$'s output $r_j(t) = y_{N-1,j}(t) \in (-1, 1)$
is used directly. For classification we sum over time,
$\hat{\ell}_j = \sum_t r_j(t)$, and feed $\hat{\ell}$ to a standard
softmax--cross-entropy loss. No learned readout weights.

\paragraph{Per task.}
\begin{itemize}[leftmargin=*]
\item \textit{Row-wise MNIST}: $N{=}28$ (one grid row per input
channel), $M{=}10$ (one bottom-shoreline cell per class). $N_{\text{ext}}
= N$, identity encoder. Decoder is summed-readout softmax-CE.
\end{itemize}

In every experiment the only trainable parameters are the per-cell
$(\log\tau_{ij}, w_{ij}^\text{L}, w_{ij}^\text{T}, b_{ij})$ quadruples.
The grid carries the whole representation.

\section{Sparse grid-RTRL via SnAP-$1$}
\label{sec:snap}

\subsection{Structural sparsity of the sensitivity tensor}

Parameter $\theta_{ij} \in \{\log\tau_{ij}, w_{ij}^\text{L},
w_{ij}^\text{T}, b_{ij}\}$ appears only in cell $(i, j)$'s update, and
its influence spreads only to downstream cells on the grid:
\begin{equation}
\frac{\partial s_{k\ell}(t)}{\partial \theta_{ij}} = 0
\quad\text{unless } k \ge i \text{ and } \ell \ge j.
\end{equation}
Unlike the one-in/one-out tree (ancestor chain of length $O(\log N)$),
the grid's two-in/two-out structure makes the downstream cone grow
quadratically with distance: at Manhattan distance $d$ from
$(i, j)$ there are $d{+}1$ descendants, and the set within distance
$k$ has size $\binom{k+2}{2}$. A full grid-RTRL would carry $O(NM)$
traces per parameter and store $\sim\!(NM)^2$ total.

\subsection{SnAP-$1$ traces}

We track self-traces for the three filter-sensitive parameters and
\emph{one cross-trace per immediate descendant} (below and right):
\begin{align}
e^{\text{L}}_{ij}(t) &\equiv \partial s_{ij}/\partial w^{\text{L}}_{ij}, &
e^{\text{T}}_{ij}(t) &\equiv \partial s_{ij}/\partial w^{\text{T}}_{ij}, &
\eta_{ij}(t) &\equiv \partial s_{ij}/\partial \log\tau_{ij}, \\
\epsilon^{(1,d)}_{ij,\theta}(t) &\equiv \partial s_d(t)/\partial \theta_{ij}, &
d &\in \{\text{below}, \text{right}\}. & &
\end{align}
Total per-cell state at SnAP-$1$: $3$ self-traces $+$ $2$ descendants
$\times\,4$ parameter types $=$ $11$ scalars, a constant factor more
than the tree primitive. Cross-trace update, for each $\theta$ and each
descendant $d$:
\begin{equation}
\epsilon^{(1,d)}_{ij,\theta}(t{+}1) =
a_d\,\epsilon^{(1,d)}_{ij,\theta}(t)
+ c_d\,\kappa_d\,(1 - y_{ij}^2)\,f_{\theta,ij}(t{+}1),
\end{equation}
where $\kappa_d$ is the weight that couples $y_{ij}$ into $d$'s drive
($w^\text{T}_d$ if $d$ is the below-neighbour, $w^\text{L}_d$ if the
right-neighbour) and $f_{\theta,ij}$ is the self-derivative.

\subsection{Backward pass}

Messages start at the bottom shoreline with
$m_{N{-}1,j}(t) = \partial L/\partial r_j(t)$ and propagate in reverse
raster order. Each cell's incoming message is the sum of contributions
from its two descendants through the spatial Jacobians
$\partial y_d/\partial y_{ij} = \kappa_d\,(1{-}y_d^2)\,c_d$. The
per-timestep gradient combines spatial + SnAP-$1$ terms with the
past-only (anti-double-count) subtraction:
\begin{equation}
\frac{\mathrm{d}L_t}{\mathrm{d}\theta_{ij}}
= m_{ij}(t)\,(1{-}y_{ij}^2)\,f_{\theta,ij}(t)
+ \!\!\sum_{d\in\{\text{below},\text{right}\}}\!\!
m_d(t)\,(1{-}y_d^2)\,\big[\epsilon^{(1,d)}_{ij,\theta}(t)
- c_d\,\kappa_d\,f_{\theta,ij}(t)\big].
\label{eq:local-grad}
\end{equation}

\section{Row-wise MNIST}
\label{sec:mnist}

Full-scale row-wise MNIST: $60{,}000$ training images, $10$ epochs,
batch $32$, Adam $\eta{=}3{\times}10^{-3}$, $\tau\in[1, 20]$\,ms,
input gain $10$. Grid dimensions match the task: $N{=}28$
(left-shoreline rows, one per input channel) and $M{=}10$
(bottom-shoreline columns, one per class). Identity encoder,
summed-readout softmax-CE decoder, no learned encoder/decoder
parameters. Local SnAP-$1$ rule (two descendant cross-traces per
cell, eleven per-cell state scalars).

\begin{table}[h]
\centering
\caption{Grid ($28{\times}10$) SnAP-$1$ on row-wise MNIST.
Single seed.}
\label{tab:mnist}
\begin{tabular}{r r r r r}
\toprule
epoch & train loss & train acc & val acc & cum.\ time \\
\midrule
0 & $2.106$ & $0.252$ & $0.374$ & $171$\,s \\
1 & $1.887$ & $0.320$ & $0.324$ & $338$\,s \\
2 & $1.783$ & $0.364$ & $0.438$ & $502$\,s \\
3 & $1.570$ & $0.467$ & $0.500$ & $665$\,s \\
4 & $1.507$ & $0.491$ & $0.503$ & $829$\,s \\
5 & $1.526$ & $0.495$ & $0.494$ & $993$\,s \\
6 & $1.584$ & $0.456$ & $0.430$ & $1155$\,s \\
7 & $1.554$ & $0.474$ & $0.479$ & $1318$\,s \\
8 & $1.476$ & $0.494$ & $0.512$ & $1485$\,s \\
9 & $1.433$ & $0.517$ & $\mathbf{0.522}$ & $1659$\,s \\
\bottomrule
\end{tabular}
\end{table}

The grid trains end-to-end on full-scale MNIST under SnAP-$1$, reaching
$0.522$ final val accuracy (vs.\ random-chance $0.10$). For reference
the tree variant at a similar per-cell state budget reports
$\sim\!0.58$ \cite{tree-neurons}; a BPTT comparison on the same grid
architecture would be welcome but is $\gg 1$\,hr per 10-epoch run on
our CPU setup and is deferred.

\paragraph{Hardware-realistic run.}
The Adam-with-batch-32 number above is a reference for the local rule
itself; for a configuration closer to what an analog chip would
actually implement we also ran plain per-group SGD with batch size
$1$---zero per-parameter optimiser state, one gradient update per
training example---at $\eta_0{=}10^{-2}$ with
$(\eta_w, \eta_b, \eta_\tau)=(\eta_0, \eta_0/20, 4\eta_0)$. Over
three epochs the validation accuracy climbs
$0.266 \to 0.283 \to 0.320$, i.e.\ $>3\times$ chance, still rising at
the end of the run. The slower convergence vs.\ Adam batch-32 is the
expected cost of dropping both variance-reduction mechanisms; the
point is that neither is \emph{required} for the local rule to drive
the architecture above chance.

\appendix

\section{Architectural ablations}
\label{app:ablations}

\paragraph{No residual feed-through.}
An earlier version of the cell included a learnable residual-feedthrough
coefficient $w_{\text{pass}}$ providing a linear shortcut around the
$\tanh$ nonlinearity. We removed it: empirical sweeps showed it offered
no material improvement in either BPTT or local-rule training after LR
ratios were tuned, and it made the forward-pass bound less clean
(residual outside $\tanh$) or killed its gradient-flow advantage
(residual inside $\tanh$). The final primitive in Section~\ref{sec:arch}
is accordingly the minimal filter+bias+$\tanh$ with no shortcut.


\begin{thebibliography}{99}

\bibitem{williams1989}
R.~J. Williams and D.~Zipser.
\newblock A learning algorithm for continually running fully recurrent neural networks.
\newblock \emph{Neural Computation}, 1(2):270--280, 1989.

\bibitem{marschall2020}
O.~Marschall, K.~Cho, and C.~Savin.
\newblock A unified framework of online learning algorithms for training recurrent neural networks.
\newblock \emph{Journal of Machine Learning Research}, 21(135):1--34, 2020.

\bibitem{tallec2018}
C.~Tallec and Y.~Ollivier.
\newblock Unbiased online recurrent optimization.
\newblock In \emph{International Conference on Learning Representations}, 2018.

\bibitem{mujika2018}
A.~Mujika, F.~Meier, and A.~Steger.
\newblock Approximating real-time recurrent learning with random Kronecker factors.
\newblock In \emph{Advances in Neural Information Processing Systems}, 2018.

\bibitem{menick2021}
J.~Menick, E.~Elsen, U.~Evci, S.~Osindero, K.~Simonyan, and A.~Graves.
\newblock A practical sparse approximation for real time recurrent learning.
\newblock In \emph{International Conference on Learning Representations}, 2021.

\bibitem{murray2019}
J.~M. Murray.
\newblock Local online learning in recurrent networks with random feedback.
\newblock \emph{eLife}, 8:e43299, 2019.

\bibitem{bellec2020}
G.~Bellec, F.~Scherr, A.~Subramoney, E.~Hajek, D.~Salaj, R.~Legenstein, and W.~Maass.
\newblock A solution to the learning dilemma for recurrent networks of spiking neurons.
\newblock \emph{Nature Communications}, 11:3625, 2020.

\bibitem{bohnstingl2023}
T.~Bohnstingl \emph{et al.}
\newblock Online spatio-temporal learning in deep neural networks.
\newblock \emph{IEEE Transactions on Neural Networks and Learning Systems}, 2023.

\bibitem{tree-neurons}
T.~Pluck.
\newblock Silicon-aware local RTRL on a hardware-primitive dendritic tree.
\newblock \url{https://github.com/ThomasPluck/tree-neurons}, 2026.

\bibitem{zucchet2023}
N.~Zucchet, R.~Meier, S.~Schug, A.~Mujika, and J.~Sacramento.
\newblock Online learning of long-range dependencies.
\newblock In \emph{Advances in Neural Information Processing Systems}, 2023.

\bibitem{orvieto2023}
A.~Orvieto, S.~L. Smith, A.~Gu, A.~Fernando, C.~Gulcehre, R.~Pascanu, and S.~De.
\newblock Resurrecting recurrent neural networks for long sequences.
\newblock arXiv:2303.06349, 2023.

\bibitem{sacramento2018}
J.~Sacramento, R.~Ponte~Costa, Y.~Bengio, and W.~Senn.
\newblock Dendritic cortical microcircuits approximate the backpropagation algorithm.
\newblock In \emph{Advances in Neural Information Processing Systems}, 2018.

\bibitem{guerguiev2017}
J.~Guerguiev, T.~P. Lillicrap, and B.~A. Richards.
\newblock Towards deep learning with segregated dendrites.
\newblock \emph{eLife}, 6:e22901, 2017.

\bibitem{payeur2021}
A.~Payeur, J.~Guerguiev, F.~Zenke, B.~A. Richards, and R.~Naud.
\newblock Burst-dependent synaptic plasticity can coordinate learning in hierarchical circuits.
\newblock \emph{Nature Neuroscience}, 24(7):1010--1019, 2021.

\bibitem{yin2024}
B.~Yin \emph{et al.}
\newblock Training high-performance deep learning networks with two-compartment dendritic neurons using e-prop.
\newblock arXiv:2402.15969, 2024.

\bibitem{young2019minatar}
K.~Young and T.~Tian.
\newblock MinAtar: An Atari-inspired testbed for thorough and reproducible reinforcement learning experiments.
\newblock arXiv:1903.03176, 2019.

\bibitem{zhang2024tclif}
S.~Zhang \emph{et al.}
\newblock TC-LIF: A two-compartment spiking neuron model for long-term sequential modelling.
\newblock arXiv:2308.13250, 2024.

\bibitem{huang2024dendsn}
X.~Huang \emph{et al.}
\newblock DendSN: A dendritic spiking neuron model for deep neural networks.
\newblock arXiv:2412.06355, 2024.

\bibitem{feng2025tslif}
L.~Feng \emph{et al.}
\newblock TS-LIF: A temporal-separation two-compartment spiking neuron.
\newblock In \emph{International Conference on Learning Representations}, 2025.

\bibitem{beniaguev2021}
D.~Beniaguev, I.~Segev, and M.~London.
\newblock Single cortical neurons as deep artificial neural networks.
\newblock \emph{Neuron}, 109(17):2727--2739, 2021.

\bibitem{iyer2022}
A.~Iyer \emph{et al.}
\newblock Avoiding catastrophe: Active dendrites enable multi-task learning in dynamic environments.
\newblock \emph{Neural Networks}, 155:1--16, 2022.

\bibitem{poirazi2020}
P.~Poirazi and A.~Papoutsi.
\newblock Illuminating dendritic function with computational models.
\newblock \emph{Nature Reviews Neuroscience}, 21(6):303--321, 2020.

\bibitem{cartiglia2022}
M.~Cartiglia, A.~Rubino, S.~Narayanan, C.~Frenkel, G.~Haessig, G.~Indiveri,
and M.~Payvand.
\newblock Stochastic dendrites enable online learning in mixed-signal neuromorphic processing systems.
\newblock arXiv:2201.10409, 2022.

\bibitem{maryada2025}
Maryada \emph{et al.}
\newblock Canonical cortical microcircuit with analog dendritic learning on neuromorphic hardware.
\newblock bioRxiv 2025.03.28, 2025.

\bibitem{acharya2021}
J.~Acharya \emph{et al.}
\newblock Dendritic computing: Branching deeper into machine learning.
\newblock \emph{Neuroscience}, 489:275--289, 2021.

\end{thebibliography}

\end{document}
